% !PRINTED 106
\weberpage{106}
%<<CONT

\[
 \varphi(h)=h f'(x)+\frac{h^2}{1\cdot2}f''(x)+\cdots+h^n a_0
\]
setzen,
\[
 |f(x)|-|\varphi(h)|<|f(x+h)|<|f(x)|+|\varphi(h)|,
\]
oder
\[
 -|\varphi(h)|<|f(x+h)|-|f(x)|<|\varphi(h)|.
\]
Nun ist, wenn $\varrho<\varepsilon$ ist, $|\varphi(h)|<\omega$, und wir können also auch sagen

\medskip
\noindent VI. \gesperrt{Der absolute Werth der Differenz}
\[
 |f(x+h)|-|f(x)|
\]
\gesperrt{ist kleiner als eine beliebig gegebene Grösse $\omega$, sobald der absolute Werth von $h$ kleiner als ein genügend kleines $\varepsilon$ ist.}
\medskip

Daraus können wir noch schliessen, wenn wir $x=y+iz$ setzen, und $h$ entweder reell oder rein imaginär annehmen, dass $|f(y+zi)|$ bei feststehendem $z$ eine stetige Function von $y$ und bei feststehendem $y$ eine stetige Function von $z$ ist.

\medskip
\noindent VII. \gesperrt{Sind die Coëfficienten $a_0,a_1,\ldots,a_n$ und die Variable $x$ reell, so kann man $x$ dem absoluten Werthe nach so gross wählen, dass das Vorzeichen von $f(x)$ mit dem Vorzeichen des ersten Gliedes $a_0x^n$ übereinstimmt, also bei positivem $a_0$ und geradem $n$ positiv, bei ungeradem $n$ für negative $x$ negativ und für positive $x$ positiv wird.}

Denn man kann in
\[
 f(x)=x^n\left(a_0+\frac{a_1}{x}+\frac{a_2}{x^2}+\cdots+\frac{a_n}{x^n}\right)
\]
$x$ dem absoluten Werthe nach so gross wählen, dass der absolute Werth von
\[
 \frac{a_1}{x}+\frac{a_2}{x^2}+\cdots+\frac{a_n}{x^n}
\]
unter den von $a_0$ heruntersinkt und also der Klammerausdruck
\[
 a_0+\frac{a_1}{x}+\frac{a_2}{x^2}+\cdots+\frac{a_n}{x^n}
\]
das Vorzeichen von $a_0$ hat, was es dann für jedes absolut grössere $x$ beibehält.
